\chapter[REFERENCIAL TEÓRICO]{REFERENCIAL TEÓRICO}
\label{cha:referencial}

Este capítulo apresenta os fundamentos teóricos que embasam o desenvolvimento deste trabalho. São abordados, em ordem narrativa: o Processo Tennessee Eastman como problema de referência; os princípios de controle plant-wide como teoria da supervisão; o diagnóstico de qualidade de malhas como o olho do supervisor; a analogia cloud-native e seus precedentes industriais; o padrão IEC~61499 e os Operators do Kubernetes como mecanismo de orquestração; o gêmeo digital como paradigma de implementação; o protocolo gRPC como canal de comunicação; o OPC UA e a família IEC~62541 como padrões de integração; e o método de integração numérica RK4 como fundamento do simulador.

% -------------------------------------------------------------------------
\section{Processo Tennessee Eastman}
\label{sec:tep}

O Processo Tennessee Eastman (\textit{Tennessee Eastman Process} — TEP) é um \textit{benchmark} amplamente utilizado na comunidade de controle de processos para avaliação de estratégias de controle, detecção de falhas e otimização de operações industriais. Proposto originalmente por Downs e Vogel em 1993, o TEP descreve uma planta química realista com dinâmica não-linear complexa, laços de reciclo e múltiplas perturbações programadas, sendo deliberadamente desafiador do ponto de vista de controle \cite{DOWNS1993}.

A planta consiste em cinco unidades de processo interconectadas: um reator, um condensador, um separador de vapor-líquido (\textit{flash separator}), um compressor e uma coluna de \textit{stripping}. Quatro reagentes (A, B, C, D) são alimentados ao reator, onde ocorrem três reações:

\begin{align}
  A + C + D &\rightarrow G \quad (\text{produto líquido}) \label{eq:rx1} \\
  A + C + E &\rightarrow H \quad (\text{produto líquido}) \label{eq:rx2} \\
  A + E     &\rightarrow F \quad (\text{subproduto gasoso}) \label{eq:rx3}
\end{align}

O componente B é inerte e não participa das reações. Os produtos G e H são os produtos desejados; F é um subproduto indesejado. As reações são exotérmicas, e a temperatura do reator é controlada indiretamente pela água de resfriamento (\textit{cooling water} — CW).

O modelo matemático é descrito por um sistema de 50 equações diferenciais ordinárias (EDO), representando os holdups de massa e energia de cada vaso de processo. O vetor de estado $\mathbf{y} \in \mathbb{R}^{50}$ inclui os inventários de cada componente nos reatores, separadores e no espaço de vapor, bem como as energias internas de cada subsistema. O modelo original foi implementado em FORTRAN e posteriormente portado para diversas linguagens, mantendo a estrutura matemática intacta.

O modelo expõe 41 variáveis medidas (XMEAS$_1$ a XMEAS$_{41}$) e 12 variáveis manipuladas (XMV$_1$ a XMV$_{12}$). As variáveis medidas incluem vazões, composições, temperaturas, pressões e níveis dos principais subsistemas; as variáveis manipuladas correspondem a posições de válvulas e setpoints de controladores de baixo nível. O ponto de operação nominal (Modo 1) é caracterizado por pressão do reator de 2705~kPa, temperatura do reator de 120,4~°C e frações molares de produto (razão G/H) específicas \cite{DOWNS1993}.

O modelo inclui 20 distúrbios programados (IDV$_1$ a IDV$_{20}$) que simulam eventos operacionais como variações de composição de \textit{feed}, falhas em trocadores de calor e contaminações. Cada distúrbio atua diretamente sobre parâmetros do modelo, permitindo avaliar a robustez de estratégias de controle frente a perturbações realistas. O modelo também define condições de parada por segurança (\textit{Interlock Shutdown Device} — ISD), que encerram a simulação quando variáveis críticas ultrapassam limites operacionais (por exemplo, pressão do reator acima de 3000~kPa ou nível do reator abaixo de 10\%).

A relevância do TEP como \textit{benchmark} decorre de sua riqueza dinâmica: o modelo combina não-linearidades, acoplamentos entre malhas de controle, laços de reciclo com elevada constante de tempo e perturbações com mecanismos físicos distintos. Essas características tornam o TEP um ambiente adequado para avaliar arquiteturas de controle supervisório — em especial aquelas que dependem de comunicação distribuída e coordenação entre componentes heterogêneos, como a arquitetura proposta neste trabalho.

% -------------------------------------------------------------------------
\section{Controle Plant-Wide}
\label{sec:plantwide}

Plantas de processo contínuo não podem ser operadas exclusivamente por malhas de controle locais. Cada controlador PID regula uma variável de seu subsistema imediato, mas é incapaz de perceber os efeitos que suas ações propagam ao longo da planta através de reciclos, transferências de massa e acoplamentos termodinâmicos. A consequência é que perturbações de grande amplitude podem atravessar toda a cadeia de processo e produzir instabilidade em pontos distantes da origem, sem que nenhum controlador local tenha autoridade para reagir \cite{Larsson2000}.

\citeonline{Larsson2000} revisamm a literatura de controle plant-wide e identificam quatro problemas centrais que qualquer estratégia de supervisão deve endereçar: (i) seleção de variáveis controladas (quais saídas manter no setpoint); (ii) seleção de variáveis manipuladas (quais entradas usar para cada objetivo); (iii) estrutura de controle (como parear variáveis); e (iv) sequenciamento de implementação (bottom-up ou top-down). A ausência de uma resposta sistemática a essas perguntas é, segundo os autores, a principal causa de estratégias de controle que funcionam individualmente mas falham quando integradas.

\citeonline{Skogestad2004} propõe uma metodologia sistemática para o projeto de estruturas de controle em plantas completas, organizada em uma hierarquia de três camadas:

\begin{enumerate}
  \item \textbf{Controle regulatório}: malhas PID locais que estabilizam variáveis de processo em torno de um ponto de operação. Esse nível opera em escala de segundos a minutos e não possui visão global da planta.
  \item \textbf{Controle supervisório}: coordenação das malhas regulatórias via ajuste de setpoints e ativação/desativação de malhas, operando em escala de minutos a horas. Esse é o nível que reage a perturbações que transcendem os limites de uma única malha local.
  \item \textbf{Otimização em tempo real}: reconfiguração de objetivos operacionais baseada em condições econômicas e restrições de processo, operando em escala de horas a dias.
\end{enumerate}

Um conceito central da metodologia de Skogestad é o de \textbf{variáveis de controle self-optimizing}: variáveis que, quando mantidas em setpoint fixo, aproximam naturalmente o ponto ótimo de operação sem necessidade de otimização explícita. A identificação dessas variáveis reduz a complexidade do supervisor, pois transforma o problema de otimização em um problema de regulação \cite{Skogestad2004}.

No contexto deste trabalho, a hierarquia de Skogestad fundamenta a separação de responsabilidades entre o \texttt{ControllerBank} (controle regulatório), o Operator Kubernetes (controle supervisório) e as extensões futuras de otimização. O TEP, com seus 41~XMEAS e 12~XMV, é um caso de estudo adequado para avaliar como um supervisor declarativo pode coordenar um conjunto de controladores regulatórios sem necessidade de reconfiguração manual.

% -------------------------------------------------------------------------
\section{Diagnóstico de Qualidade de Malhas}
\label{sec:diagnostico}

A qualidade de uma malha de controle não é binária: um controlador pode estar operando, mas de forma subótima. Malhas com ganho excessivo oscilam; malhas com ganho insuficiente respondem lentamente e permitem que distúrbios se propaguem; malhas com setpoint desatualizado operam em ponto de trabalho inadequado. Em instalações industriais com centenas ou milhares de malhas, identificar automaticamente quais controladores estão mal sintonizados é um problema relevante e de difícil solução manual \cite{Bradu2018}.

\subsection{Índice de Harris e o Benchmark de Variância Mínima}

O trabalho seminal de \citeonline{Harris1989} estabeleceu o primeiro índice quantitativo para avaliação automática de desempenho de malhas fechadas. A ideia central é que, para qualquer processo linear com distúrbios estocásticos, existe um limite inferior para a variância da variável de processo: a variância do controlador de variância mínima. Esse limite pode ser estimado diretamente dos dados de operação rotineira, sem necessidade de identificação explícita do modelo do processo.

O índice de Harris $\eta \in [0, 1]$ compara a variância observada $\sigma^2_y$ com o limite de variância mínima $\sigma^2_{\text{mv}}$:

\begin{equation}
  \eta = \frac{\sigma^2_{\text{mv}}}{\sigma^2_y}
  \label{eq:harris_index}
\end{equation}

Valores próximos de 1 indicam operação próxima ao ótimo estocástico; valores baixos indicam desempenho degradado. A estimativa de $\sigma^2_{\text{mv}}$ é obtida ajustando-se um modelo autorregressivo de ordem $d$ (atraso puro do processo) aos dados de operação, sem necessidade de experimentos de identificação \cite{Harris1989}. O índice tornou-se referência na literatura de \textit{control loop performance monitoring} (CLPM) e motivou uma geração de índices derivados, incluindo o Índice de Preditividade descrito a seguir.

\subsection{Índice de Preditividade e a Abordagem do CERN}

O sistema criogênico do LHC opera aproximadamente 5.000 malhas PID para manter os imãs supercondutores a 1,9~K. Nessa escala, qualquer abordagem de diagnóstico baseada em inspeção humana é inviável: a revisão manual de 5.000 malhas levaria centenas de horas de trabalho especializado por ciclo de operação \cite{Bradu2018}.

\citeonline{Bradu2018} propõem o \textbf{Índice de Preditividade} (\textit{Predictability Index} — PI) como métrica automática para triagem de malhas. O PI mede a capacidade de um modelo autorregressivo de prever o comportamento futuro da variável de processo: malhas bem reguladas são inerentemente previsíveis (pequenos resíduos de predição), enquanto malhas oscilantes ou instáveis apresentam comportamento imprevisível (resíduos elevados). Formalmente, o PI é calculado como:

\begin{equation}
  \text{PI} = 1 - \frac{\text{RMSE}(\hat{y}, y)}{\sigma_y}
  \label{eq:pi_index}
\end{equation}

\noindent onde $\hat{y}$ é a predição do modelo autorregressivo de ordem $m$ com horizonte $b$, $y$ é o valor real e $\sigma_y$ é o desvio padrão da série. O PI varia de 0 (completamente imprevisível) a 1 (perfeitamente previsível); um limiar prático para diagnóstico é tipicamente fixado em PI $< 0{,}7$ para sinalizar malhas problemáticas \cite{Bradu2018}.

Trabalhos complementares do mesmo grupo investigam a detecção de oscilações via métodos baseados em conhecimento especializado \cite{Tilaro2017} e o aprendizado de modelos para detecção de anomalias em sistemas de controle do CERN \cite{Tilaro2018}. Em conjunto, esses trabalhos estabelecem que o diagnóstico automático de malhas em larga escala é viável com técnicas baseadas em análise de séries temporais, sem necessidade de modelos físicos explícitos.

\subsection{Relevância para este Trabalho}

O gêmeo digital desenvolvido neste trabalho é capaz de gerar, de forma controlada e reprodutível, séries temporais das 41~XMEAS sob qualquer condição de operação — com controladores bem sintonizados, degradados, e sob distúrbios IDV específicos. Essa capacidade de experimentação controlada é o que torna o laboratório TEP uma plataforma natural para cálculo do Índice de Preditividade: os históricos sintéticos podem ser usados para calibrar os parâmetros $(m, b, n)$ do PI, identificar os limiares adequados para cada variável, e validar a lógica de diagnóstico antes de aplicá-la a dados de plantas reais.

% -------------------------------------------------------------------------
\section{Cloud-Native como Paradigma de Orquestração Industrial}
\label{sec:cloudnative}

O termo \textit{cloud-native} designa uma abordagem de desenvolvimento de sistemas que explora as capacidades da computação em nuvem para construir e executar aplicações escaláveis, resilientes e gerenciáveis. Os princípios fundamentais — conteinerização, microsserviços, orquestração declarativa e entrega contínua — foram consolidados a partir da experiência acumulada por empresas de tecnologia de grande porte ao gerenciar infraestruturas de escala global \cite{Shuiguang2023}.

\citeonline{Shuiguang2023} apresentam uma revisão abrangente da computação cloud-native sob a perspectiva de serviços, identificando que a essência do paradigma é a separação entre a definição do que um sistema deve fazer (declaração de estado desejado) e a mecânica de como esse estado é alcançado e mantido (infraestrutura de orquestração). Essa separação é análoga ao problema central do controle supervisório: o operador declara uma política de operação, e a infraestrutura de controle trabalha para manter a planta dentro dessa política.

\citeonline{Burns2016} documentam a evolução dos sistemas de orquestração do Google — Borg, Omega e Kubernetes — e identificam que o salto qualitativo de Borg para Kubernetes foi exatamente a adoção de uma API declarativa baseada em estado desejado, em contraste com as APIs imperativas anteriores. Esse princípio, que o Google consolidou ao longo de mais de uma década operando clusters com dezenas de milhares de máquinas, é o mesmo que fundamenta a proposta de usar o Kubernetes como plataforma supervisória neste trabalho.

\subsection{Precedentes Industriais}

A aplicação de princípios cloud-native a sistemas de controle industrial não é apenas uma analogia: trabalhos recentes demonstram implementações concretas.

\citeonline{Johansson2022} demonstram a viabilidade de usar o Kubernetes para orquestrar sistemas de controle distribuído de alta disponibilidade (\textit{High Availability Distributed Control Systems}). Os autores constroem um sistema em que os componentes de controle são implantados como pods Kubernetes e o plano de controle nativo gerencia failover, reinicialização e distribuição de carga — funções que em sistemas de controle industriais tradicionais exigem hardware redundante dedicado. Os resultados demonstram que o Kubernetes pode fornecer os mesmos níveis de disponibilidade com infraestrutura de propósito geral.

\citeonline{Mellado2022} propõem o \textbf{IoT-PLC}, um CLP conteinerizado para a Indústria~4.0 em que cada funcionalidade — controle regulatório, armazenamento de dados de campo, interfaces sem fio — é encapsulada em um contêiner separado. O IoT-PLC usa um \textit{modelo de dispositivo virtual} que abstrai a interface com camadas superiores, permitindo que aplicações na nuvem interajam com o controlador físico de forma transparente. Esse modelo de abstração é precursor direto do CRD \texttt{PLCMachine} proposto neste trabalho.

A convergência desses trabalhos sugere que a integração entre infraestrutura cloud-native e sistemas de controle industrial não é uma extrapolação especulativa, mas uma direção ativa de pesquisa com resultados publicados e reproduzíveis.

% -------------------------------------------------------------------------
\section{IEC~61499 e Controle Distribuído por Blocos de Função}
\label{sec:iec61499}

A norma \textbf{IEC~61499} \cite{IEC61499} define uma arquitetura para sistemas de controle distribuído baseada em \textbf{Blocos de Função} (\textit{Function Blocks} — FBs). Proposta como evolução da IEC~61131-3 (que standardiza a programação de CLPs individuais), a IEC~61499 endereça o problema da distribuição: como decompor uma aplicação de controle em unidades modulares que possam ser implantadas em dispositivos distintos e comunicar-se entre si via redes industriais.

Um Bloco de Função na IEC~61499 encapsula tanto dados (variáveis de entrada/saída) quanto comportamento (algoritmos e máquinas de estados de eventos). Os FBs se comunicam por \textbf{eventos} (que disparam execução) e \textbf{dados} (que carregam valores de processo). Essa separação entre fluxo de controle (eventos) e fluxo de dados é análoga à separação entre plano de controle e plano de dados em redes de computadores — o que facilita a distribuição dos componentes por múltiplos dispositivos físicos.

A norma define dois níveis de abstração:

\begin{enumerate}
  \item \textbf{Modelo de aplicação}: a lógica de controle descrita como um grafo de FBs interconectados, independente de hardware.
  \item \textbf{Modelo de sistema}: a distribuição dos FBs por dispositivos físicos e os recursos (\textit{resources}) dentro de cada dispositivo que executam os FBs.
\end{enumerate}

A distinção entre modelo de aplicação e modelo de sistema é o elemento da IEC~61499 mais relevante para este trabalho: o Operator Kubernetes realiza exatamente a função de um \textit{deployment manager} da IEC~61499 — mapeia a aplicação (lógica supervisória declarada no CRD) para o sistema (componentes em execução no cluster). A \textit{spec} do CRD corresponde ao modelo de aplicação; o \textit{status} corresponde ao estado observado do sistema; o \textit{controller loop} realiza a reconciliação.

No Capítulo~5, a correspondência entre o padrão de Operator e os princípios da IEC~61499 é analisada em detalhes, concluindo que o Kubernetes pode servir como infraestrutura de runtime para aplicações modeladas segundo a norma, desde que a lógica de cada FB seja encapsulada em contêineres com interface gRPC.

% -------------------------------------------------------------------------
\section{Kubernetes e o Padrão de Operators}
\label{sec:kubernetes}

O Kubernetes é uma plataforma de orquestração de contêineres de código aberto desenvolvida originalmente pelo Google e atualmente mantida pela \textit{Cloud Native Computing Foundation} (CNCF). Sua arquitetura é centrada no princípio da \textbf{reconciliação declarativa}: o usuário descreve o estado desejado do sistema em manifestos YAML, e o plano de controle do Kubernetes trabalha continuamente para fazer o estado real convergir ao estado desejado \cite{Burns2016}.

A arquitetura do Kubernetes divide-se em plano de controle (\textit{control plane}) e nós de trabalho (\textit{worker nodes}). O plano de controle inclui o \textit{API Server} (único ponto de entrada para todas as operações), o \textit{etcd} (banco de dados distribuído de chave-valor que armazena o estado do cluster), o \textit{Scheduler} (alocador de cargas de trabalho) e o \textit{Controller Manager} (gerenciador dos controladores nativos). Os nós de trabalho executam os contêineres via \textit{kubelet} e \textit{container runtime}.

O Kubernetes permite estender seu modelo de recursos por meio das \textbf{Definições de Recursos Customizados} (\textit{Custom Resource Definitions} — CRDs). Um CRD declara um novo tipo de recurso no API Server, com seu próprio esquema (\textit{spec} e \textit{status}), tornando-o disponível para criação, listagem e edição com as mesmas ferramentas (kubectl, YAML) usadas para recursos nativos. A separação entre \textit{spec} (intenção declarada pelo usuário) e \textit{status} (estado observado pelo sistema) é fundamental para o padrão de Operator \cite{Burns2016}.

O \textbf{padrão de Operator} (\textit{Operator Pattern}) é um método de extensão do Kubernetes que codifica o conhecimento operacional de uma aplicação específica em um controlador customizado \cite{Burns2016}. Um Operator combina um ou mais CRDs com um \textit{controller loop} que observa instâncias desses recursos e age para reconciliar o estado real com o estado desejado — substituindo a intervenção humana pela automação do domínio de conhecimento específico da aplicação.

O \textit{loop} de reconciliação de um Operator segue o padrão \textbf{Observar–Avaliar–Agir}:
\begin{enumerate}
  \item \textbf{Observar}: o controller recebe uma notificação de que um recurso customizado foi criado, modificado ou deletado (via \textit{watch} no API Server).
  \item \textbf{Avaliar}: o controller lê o estado atual da aplicação (possivelmente via chamadas externas, como gRPC) e compara com o estado desejado descrito no \textit{spec} do CRD.
  \item \textbf{Agir}: se houver divergência, o controller toma ações corretivas — no caso deste trabalho, chamadas gRPC à planta para ajustar parâmetros de controle.
\end{enumerate}

O framework \textbf{Kubebuilder} é a ferramenta oficial da CNCF para construção de Operators em Go. Ele provê scaffolding de projeto, integração com o \textit{controller-runtime} (biblioteca Go para interação com o API Server), geração de CRDs a partir de structs Go anotadas, e utilitários para testes.

Para operações industriais, o padrão de Operator oferece vantagens relevantes: a lógica supervisória é codificada uma vez e executada de forma autônoma; a política de operação é declarativa e auditável em repositórios de versionamento; e o Kubernetes garante alta disponibilidade do próprio operator via restarts automáticos \cite{Burns2016}.

% -------------------------------------------------------------------------
\section{Gêmeo Digital Industrial}
\label{sec:twin}

\subsection{Conceito Fundador e Origem}

O conceito de \textbf{gêmeo digital} (\textit{digital twin}) origina-se do programa Apollo da NASA, onde gêmeos idênticos de naves espaciais eram mantidos na Terra para espelhar as condições das naves em órbita, permitindo que engenheiros simulassem alternativas e assistissem os astronautas em situações críticas \cite{Boschert2016}. A ideia foi posteriormente adaptada para sistemas terrestres: na indústria aeroespacial, o conceito do \textit{Iron Bird} — uma integração física de sistemas elétricos, hidráulicos e controles de voo com layouts correspondentes aos aviões reais — representa um análogo físico-digital que permite otimização e validação de falhas em estágios iniciais de desenvolvimento \cite{Boschert2016}.

Com o avanço das tecnologias de simulação e integração de dados, o conceito evoluiu para modelos totalmente digitais que acompanham o ativo físico ao longo de todo seu ciclo de vida. O \textbf{gêmeo digital industrial} moderno refere-se a uma \textit{representação digital compreensiva e funcional de um componente, produto ou sistema que inclui informações úteis em todas as fases do ciclo de vida} \cite{Boschert2016}.

\subsection{Gêmeo Digital e Sistemas Cyber-Físicos}

O conceito de \textbf{Sistema Cyber-Físico} (\textit{Cyber-Physical System} — CPS) designa sistemas em que componentes computacionais e físicos estão profundamente integrados, comunicando-se continuamente para que o mundo digital possa influenciar o mundo físico e vice-versa \cite{Hehenberger2016}. O CPS representa o substrato conceitual sobre o qual o gêmeo digital opera: enquanto o gêmeo digital é a representação, o CPS é a arquitetura que mantém a sincronização entre representação e realidade.

\citeonline{Hehenberger2016} discutem a evolução dos sistemas mecatrônicos em direção a CPSs, identificando que o salto qualitativo não é apenas de complexidade, mas de paradigma: em um CPS, a fronteira entre o artefato físico e seu modelo digital é deliberadamente porosa, e a lógica de controle pode residir em qualquer lado dessa fronteira. Essa perspectiva é relevante para este trabalho: o laboratório TEP é um CPS prototípico em que a planta física (simulada em Rust) e o supervisor (Operator Kubernetes) coexistem em camadas distintas mas comunicantes via gRPC.

\subsection{Níveis de Fidelidade}

Uma taxonomia comum classifica os gêmeos digitais em três níveis de fidelidade \cite{Boschert2016}:
\begin{enumerate}
  \item \textbf{Modelo digital}: representação estática ou off-line, sem sincronização em tempo real com o ativo.
  \item \textbf{\textit{Digital shadow}}: representação que recebe dados do ativo físico mas não envia comandos de volta; usada para monitoramento e diagnóstico.
  \item \textbf{Gêmeo digital completo}: integração bidirecional — recebe dados do ativo e pode enviar atuações e comandos, formando um \textit{cyber-physical loop} que fecha sobre operação e controle.
\end{enumerate}

O laboratório desenvolvido neste trabalho implementa um gêmeo digital do tipo \textit{shadow} em transição para \textit{completo}: a planta simulada em Rust replica a dinâmica do TEP via modelo matemático de 50 EDOs, e o Operator Kubernetes atua sobre ela em um laço de controle supervisório. Embora não haja uma planta física correspondente no presente trabalho, a arquitetura é deliberadamente construída segundo o paradigma de gêmeo digital: a substituição da planta simulada por dados reais requeriria apenas a troca do endereço gRPC de destino, preservando toda a lógica supervisória.

\subsection{Limite Honesto da Implementação}

É importante delimitar o que foi entregue neste trabalho em relação ao espectro completo de um gêmeo digital. O laboratório TEP implementa o modelo matemático original com fidelidade, mas não inclui: modelos de degradação de equipamentos, estimação de estado baseada em dados reais, ou sincronização com um ativo físico. O sistema é mais precisamente classificado como um \textbf{simulador com interface supervisória} — um passo necessário antes de um gêmeo digital operacional, mas não equivalente a ele \cite{Boschert2016}.

% -------------------------------------------------------------------------
\section{gRPC como Protocolo de Comunicação}
\label{sec:grpc}

O gRPC é um \textit{framework} de chamada de procedimento remoto (\textit{Remote Procedure Call} — RPC) de alto desempenho desenvolvido pelo Google, baseado no protocolo HTTP/2 e no sistema de serialização \textit{Protocol Buffers} (protobuf) \cite{GRPC}. Diferentemente de alternativas baseadas em REST/JSON, o gRPC utiliza serialização binária tipada, o que reduz o tamanho das mensagens e o tempo de serialização — vantagem relevante para aplicações de automação que transmitem dados de sensores em alta frequência.

A definição de um serviço gRPC é feita em um arquivo \texttt{.proto}, que especifica os métodos RPC e os tipos de mensagem de forma agnóstica à linguagem. A partir desse arquivo, o compilador \texttt{protoc} gera \textit{stubs} de cliente e servidor para as linguagens desejadas — no caso deste trabalho, Rust (biblioteca \textit{tonic}) para a planta, Go (\textit{google.golang.org/grpc}) para o operator, e Python (\textit{grpcio}) para a IHM.

O gRPC suporta quatro modalidades de comunicação:
\begin{enumerate}
  \item \textbf{Unário}: uma requisição e uma resposta — usado para comandos pontuais (\texttt{GetPlantStatus}, \texttt{UpdateController}).
  \item \textbf{\textit{Server streaming}}: uma requisição e um fluxo contínuo de respostas — usado para streaming de métricas (\texttt{StreamMetrics}).
  \item \textbf{\textit{Client streaming}}: fluxo de requisições e uma resposta.
  \item \textbf{Bidirecional}: fluxos simultâneos em ambas as direções.
\end{enumerate}

No contexto industrial, o \textit{server streaming} é particularmente adequado para a transmissão contínua de dados de processo: a IHM abre um único fluxo gRPC e recebe atualizações das 41~XMEAS e 12~XMV a cada passo de simulação, sem \textit{polling} periódico. O protocolo HTTP/2 subjacente permite múltiplos fluxos simultâneos em uma única conexão TCP, reduzindo overhead de estabelecimento de conexão.

% -------------------------------------------------------------------------
\section{OPC UA e a Família de Normas IEC~62541}
\label{sec:opcua}

\subsection{Contexto e Motivação}

Antes do OPC UA, a comunicação entre sistemas de automação industrial era dominada por protocolos proprietários e incompatíveis. O OPC clássico (OPC DA, OPC AE, OPC HDA), lançado em 1996, foi a primeira tentativa de padronização — mas era baseado em COM/DCOM da Microsoft, o que o restringia a Windows e criava barreiras de segurança e escalabilidade.

O \textbf{OPC UA} (\textit{OPC Unified Architecture}), desenvolvido pela OPC Foundation a partir de 2006 e formalizado como família de normas internacionais \textbf{IEC~62541} \cite{IEC62541_1}, representa uma ruptura com esse modelo: é independente de plataforma e sistema operacional, orientado a serviços, com modelo de segurança nativo baseado em criptografia de chave pública (X.509) e transportável sobre qualquer rede IP \cite{Mahnke2009intro}. A norma unifica os padrões clássicos — acesso a dados em tempo real (DA), alarmes e eventos (AE) e acesso histórico (HDA) — em uma arquitetura coesa e extensível.

\subsection{Arquitetura}

A arquitetura do OPC UA é centrada no conceito de \textbf{espaço de endereçamento} (\textit{Address Space}): um grafo de nós (\textit{Nodes}) interconectados por referências tipadas (\textit{References}), que representa de forma uniforme todos os dados, metadados, métodos e eventos expostos por um servidor \cite{Mahnke2009arch}. Cada nó possui uma classe (\textit{NodeClass}) que define seu papel:

\begin{itemize}
  \item \textbf{Object}: representa entidades do mundo real (e.g., um reator, um sensor).
  \item \textbf{Variable}: armazena um valor escalar ou estruturado com timestamp, qualidade e tipo de dado.
  \item \textbf{Method}: define uma operação executável (e.g., ligar uma bomba).
  \item \textbf{ObjectType / VariableType}: tipos reutilizáveis que definem modelos de informação.
\end{itemize}

Um cliente OPC UA navega o espaço de endereçamento usando o serviço \texttt{Browse}, lê e escreve valores com \texttt{Read}/\texttt{Write}, e recebe notificações de mudança via \texttt{Subscription}/\texttt{MonitoredItem} \cite{Mahnke2009services}. O servidor publica atualizações periodicamente apenas quando o valor monitorado muda além de um limiar configurável (\textit{DeadBand}), reduzindo tráfego de rede desnecessário.

\subsection{Família de Normas IEC~62541}

A IEC~62541 é composta por 14 partes, cada uma cobrindo um aspecto distinto da especificação. A Tabela~\ref{tab:iec62541} apresenta as partes mais relevantes para este trabalho.

\begin{table}[ht!]
\centering
\caption{Partes selecionadas da família IEC~62541 e sua relevância para o projeto TEP}
\label{tab:iec62541}
\begin{tabular}{|c|l|p{5cm}|p{3.2cm}|}
\hline
\textbf{Parte} & \textbf{Título} & \textbf{Versa sobre} & \textbf{Edições} \\
\hline
1 & Visão Geral e Conceitos & Fundamentos, terminologia e motivação do padrão & 2009, 2025 \\
\hline
3 & Modelo do Espaço de End. & Nós, referências, navegação — base de todo modelo de informação \cite{IEC62541_3} & 2010, 2020, 2025 \\
\hline
4 & Serviços & API completa do servidor: Read, Write, Browse, Subscribe, Call & 2010, 2020 \\
\hline
8 & Acesso a Dados & Variáveis analógicas com unidade de engenharia, faixa e qualidade — expor XMEAS/XMV \cite{IEC62541_8} & 2011, 2020, 2025 \\
\hline
9 & Alarmes e Condições & Modelo de alarmes: estados, reconhecimento, supressão — mapeável aos ISDs do TEP & 2012, 2020 \\
\hline
11 & Acesso Histórico & Séries temporais: leitura e agregação de dados históricos & 2015, 2020 \\
\hline
14 & PubSub & Publish-subscribe sobre MQTT/AMQP/UDP — complemento IIoT & 2020, 2026 \\
\hline
\end{tabular}
\fonte{Elaborado pelo autor com base em \citeonline{IEC62541_1} e \citeonline{Mahnke2009intro}.}
\end{table}

A \textbf{Parte~8 (Acesso a Dados)} \cite{IEC62541_8} é a mais relevante para expor variáveis de processo: o tipo \texttt{AnalogItemType} encapsula um valor numérico com sua unidade de engenharia, faixa operacional e faixa instrumental — metadados que as XMEAS do TEP possuem intrinsecamente (XMEAS$_7$, por exemplo, é pressão em kPa com faixa operacional definida). A qualidade (\textit{Bad}/\textit{Uncertain}/\textit{Good}) associada a cada leitura permite que clientes OPC UA tomem decisões com base na confiabilidade dos dados.

\subsection{Integração Empresarial e o Modelo Purdue — IEC~62264}

A norma \textbf{IEC~62264} \cite{IEC62264} (\textit{Enterprise-Control System Integration}) define modelos e terminologia para integração entre sistemas de gestão empresarial (Nível~4 — ERP) e sistemas de controle de manufatura (Níveis~0 a~3). O modelo de referência hierárquico da IEC~62264 — derivado do Modelo Purdue — organiza as funções de controle em camadas com responsabilidades e horizontes temporais distintos, alinhando-se à hierarquia de Skogestad descrita na Seção~\ref{sec:plantwide}.

\subsection{OPC UA no Contexto deste Trabalho}

A comunicação interna do laboratório foi implementada com gRPC por suas vantagens de desempenho e integração nativa com Rust, Go e Python. No entanto, o gRPC é um protocolo interno — não é reconhecido por sistemas SCADA, CLPs ou ferramentas de terceiros. O OPC UA, por sua vez, é o protocolo de interoperabilidade industrial adotado por praticamente todos os fabricantes relevantes e normalizado pela IEC~62541.

A integração futura prevista envolve expor a planta TEP como um servidor OPC UA, mapeando as 41~XMEAS como variáveis \texttt{AnalogItemType} (Parte~8) e os distúrbios IDV como métodos invocáveis (Parte~4). Isso permitiria: (a) conectar controladores reais ao gêmeo digital via OPC UA, substituindo o \texttt{ControllerBank} interno; (b) integrar ferramentas de diagnóstico baseadas em OPC UA para cálculo do Índice de Preditividade; e (c) arquivar os dados de experimentos como séries históricas acessíveis via Parte~11.

% -------------------------------------------------------------------------
\section{Integração Numérica por Runge-Kutta de 4ª Ordem}
\label{sec:rk4}

A simulação dinâmica do TEP requer a integração numérica de um sistema de 50~EDOs da forma:

\begin{equation}
  \frac{d\mathbf{y}}{dt} = f(\mathbf{y}, \mathbf{u}, t)
  \label{eq:ode}
\end{equation}

\noindent onde $\mathbf{y} \in \mathbb{R}^{50}$ é o vetor de estado, $\mathbf{u} \in \mathbb{R}^{12}$ é o vetor de variáveis manipuladas e $f$ é a função de dinâmica da planta — altamente não-linear, com acoplamentos termodinâmicos e de transferência de massa.

O método de \textbf{Runge-Kutta de 4ª ordem} (RK4) é o integrador clássico para EDOs de valores iniciais com dinâmica suave. Para um passo de tempo $h$, a atualização do estado é:

\begin{align}
  k_1 &= h \cdot f(\mathbf{y}_n, t_n) \nonumber \\
  k_2 &= h \cdot f\!\left(\mathbf{y}_n + \tfrac{k_1}{2},\; t_n + \tfrac{h}{2}\right) \nonumber \\
  k_3 &= h \cdot f\!\left(\mathbf{y}_n + \tfrac{k_2}{2},\; t_n + \tfrac{h}{2}\right) \nonumber \\
  k_4 &= h \cdot f(\mathbf{y}_n + k_3,\; t_n + h) \nonumber \\
  \mathbf{y}_{n+1} &= \mathbf{y}_n + \tfrac{1}{6}(k_1 + 2k_2 + 2k_3 + k_4)
  \label{eq:rk4}
\end{align}

O RK4 possui erro local de truncamento de ordem $\mathcal{O}(h^5)$ e erro global de ordem $\mathcal{O}(h^4)$. Para o TEP, o modelo original especifica $dt = 0{,}001$~h como passo de integração padrão, valor que garante estabilidade numérica para a dinâmica mais rápida do sistema \cite{DOWNS1993}.

A norma dos derivativos $\|\mathbf{f}(\mathbf{y}, t)\|_\infty$ — denominada \texttt{deriv\_norm} na implementação — serve como indicador de saúde numérica da simulação: valores excessivamente altos indicam instabilidade iminente; valores estabilizados em patamar baixo indicam operação próxima ao estado estacionário.

% -------------------------------------------------------------------------
\section{Considerações Finais}
\label{sec:ref_conclusao}

A revisão apresentada fundamenta as escolhas tecnológicas do trabalho segundo a cadeia argumentativa do Big Picture: o TEP provê o problema que força a pergunta sobre supervisão; o controle plant-wide (Skogestad, Larsson) fornece a teoria da resposta; o diagnóstico de malhas (Harris, Bradu) define o que o supervisor deve observar; a analogia cloud-native (Shuiguang, Burns) e os precedentes industriais (Johansson, Mellado) mostram que a abordagem é viável; o IEC~61499 formaliza o modelo de dispositivo sobre o qual o Operator opera; o gêmeo digital enquadra o laboratório como arquitetura válida para investigação de orquestração industrial; e o gRPC e o OPC UA (IEC~62541) formalizam as interfaces entre os subsistemas. Os capítulos seguintes descrevem como esses elementos foram combinados na prova de conceito.
